\section{The Propers: Detailed Commentary}

\subsection{Introit: the promise to the exiles}

Jeremiah's letter addresses people carried away from Jerusalem to Babylon.
Its instructions are unusually practical: build houses, plant gardens,
form households, and pray for the peace of the city in which they live.
Their life there will last long enough for such work to matter. The false
promise of a speedy return would make those responsibilities seem
pointless. The Lord's promised peace instead permits a faithful life during
displacement.\footnote{Jer 29:1--14; compare the dispute over an early
return in Jer 28. Biblical context follows the registered
Douay--Rheims--Challoner text.}

The chant combines parts of three verses. It omits the whole-hearted
seeking of verse 13 and the fuller description of gathering in verse 14.
Its brevity concentrates on God's initiative, invocation and answer.
The omitted context remains necessary for understanding whose peace is
promised and why the people must wait. Jerome explains the promise against
the false prophets who announce an immediate return. He also asks why
people must pray for a return God has already promised. His answer
concerns their readiness to receive it. That argument will receive its
full treatment in the second interpretation.\footnote{Jerome,
\work{In Ieremiam} V, on 29:10--14, PL 24, cols. 893--896.}

The attached psalm verse puts restoration in the past tense, while the
complete Psalm 84(85) goes on to ask for conversion and renewed life.
Augustine explains the apparent reversal through prophecy: the speaker
announces as accomplished what remains future from his standpoint.
The subsequent petition for conversion shows that the announced restoration
is still awaited. Augustine also understands captivity as bondage to sin
and death, and follows the Psalm through to its promised peace and fruitful
land.\footnote{Augustine, \work{Enarrationes in Psalmos} 84.1--13,
especially §§2--4; NPNF English heading, Psalm LXXXV. The complete Psalm
is Ps 84(85).}

\subsection{Collect: wills capable of bearing fruit}

The Collect addresses God on behalf of his faithful. Their membership in
the Church is already acknowledged when the prayer asks him to arouse
their wills. Its request includes their active pursuit of fruit:
grace makes willing and doing possible. The further request for
\textit{remedia}, remedies, concerns people already pursuing the divine
work. The prayer describes growth in dependence on God's goodness.

Lasance's historical English renders the medical noun as ``mercies.''
That broader rendering should be read beside the Latin, whose reference
to remedies anticipates the Postcommunion's explicit cure. Lucien Fromage,
in the continuation of \work{The Liturgical Year}, explains the Collect
through grace-assisted works preparing for greater grace. He attributes
the movement of the will to God. His comment concerns this prayer
directly, whereas Aquinas's explanation of fruitful knowledge in the
Epistle concerns a different text.\footnote{Fromage, continuation of
\work{The Liturgical Year}, English trans. Laurence Shepherd, vol. XI
(Dublin: Duffy, 1883), pp. 503--504; Aquinas,
\work{Super Colossenses} 1, lectio 3.}

\subsection{Epistle: knowledge expressed in a life}

Paul has heard of the Colossians' faith and love before he prays for their
further growth. His request concerns knowledge of God's will, conduct
pleasing to God, fruitful work, endurance and thanksgiving. These are
requests for people already addressed as believers. The whole letter
continues with the Son's relation to creation, reconciliation, warnings
against deception, and the practices of a renewed life. Its later command
to forgive belongs to that renewed life.\footnote{Col 1:2--23;
2:1--10; 3:1--17, especially 3:13.}

The appointed introduction replaces the continuing letter's opening.
The Missal also prints ``the will of God'' where the continuous
Clementine and its Challoner English have ``his will.'' The reference to
redemption through Christ's Blood belongs to the appointed Latin
tradition. These details identify the actual lesson; a Bible passage at
the same address is not automatically its exact liturgical wording.

Chrysostom explains the repeated requests for fullness as both
encouragement and correction. The recipients have received grace,
but still lack its full fruit. Paul can praise what exists without
encouraging them to stop. Aquinas observes a further progression:
knowledge should produce action, and faithful action disposes a person
to deeper knowledge. Growth includes endurance when action becomes
difficult. Their detailed accounts of rescue and worthiness will
distinguish the two interpretations below.\footnote{Chrysostom,
\work{Homilies on Colossians} 2, on 1:9--14; Aquinas,
\work{Super Colossenses} 1, lectio 3.}

\subsection{Gradual: praise within a psalm of affliction}

The whole Psalm 43(44) recalls what the ancestors received through God's
action. Possession of the land was not achieved by their sword or by
confidence in their own strength. The appointed verses praise God for
rescue and promise enduring thanksgiving. The Psalm then describes
defeat, scattering and humiliation, although the speakers insist that
they have remained faithful. Its latter half prevents the chant from
being used as a promise that faithful people always win their battles.

Augustine and Bellarmine answer the Psalm's change of tense differently.
Augustine hears future deliverance announced with the certainty of an
accomplished fact. Bellarmine hears praise for actual past rescues before
the lament over present suffering. Their explanations of the praise
cannot be substituted for one another, although both permit an afflicted
people to hope. The first interpretation develops the reasons for that
difference.\footnote{Augustine, \work{Enarrationes} 43.5--8, NPNF heading
XLIV; Bellarmine, \work{Commentary on the Book of Psalms}, Psalm 43,
explanation of native vv. 7--9, in John O'Sullivan's abridged English.
These native verse comments concern the Mass's Ps 43:8--9.}

\subsection{Alleluia: a prayer from the depths}

The Alleluia appoints only the opening cry of Psalm 129(130). Its
liturgical Latin asks God to hear the speaker's prayer, while the
continuous Psalm asks him to hear the voice and continues with an
appeal to his attentive ears. The chant's shorter wording is complete
as a chant.

Augustine begins his exposition with Jonah praying beneath the sea and
within the fish. No bodily depth prevents God from hearing. He then
interprets the depths as mortal life and the burden of sin. The decisive
distinction is between someone who cries out and someone who has become
content with the depths. The former still seeks deliverance. The latter
can even mistake prosperity in sin for proof that nothing is wrong.
The Alleluias therefore accompany a genuine plea, rather than announce
that the speaker has ceased to need mercy.\footnote{Augustine,
\work{Enarrationes} 129.1--2; NPNF heading CXXX.}

\subsection{Gospel: flight, deception and gathering}

Jesus speaks after leaving the Temple. Seated on Mount Olivet, he
answers questions about the buildings' destruction, his coming and the
consummation. The appointed passage begins within that answer.
Its reference to Daniel's abomination of desolation recalls the
prophet's descriptions of desecration and persecution. Daniel also
describes the Son of man coming with the clouds and receiving dominion.
Both the threatened holy place and the coming ruler have scriptural
antecedents.\footnote{Mt 24:1--3, 15--35; Dan 7:12--14; 9:20--27;
11:31; 12:1--13. These Daniel passages are explanatory intertexts,
not additional lessons appointed for this Mass.}

The injunction to flee concerns people in Judea. The roof, field,
pregnancy, nursing, winter and Sabbath make the danger concrete.
Chrysostom explains these details in relation to escape from the
catastrophe. They deserve to be understood before a spiritual
application is made. Jesus then warns against a Christ located by
rumour in a desert or a hidden room. His own coming will be manifest
like lightning, followed by gathering from the four winds.

The passage ends with a fig tree, a statement about ``this generation,''
and Christ's abiding words. A seasonal sign can assure the hearer
that something is approaching without naming its day. The first
interpretation follows Chrysostom's explanation of the transition
between the warnings and of the generation that survives them.
The following verse's unknown day and hour, and the calls to
watchfulness that follow, remain part of the discourse's
context.\footnote{Chrysostom, \work{Homilies on Matthew} 75--77;
Mt 24:36--51 and the continuing judgment teaching in Mt 25.}

\subsection{Offertory: the cry repeated}

The Offertory returns to the same Psalm as the Alleluia, but repeats its
opening after the petition to be heard. Its exact words and ritual place
are distinct. The second cry accompanies the offering rather than the
approach to the Gospel.

The repetition allows the worshipper to continue asking after hearing
Christ's warnings. Connected with the Secret, it can also be understood
as a request that the offerer receive conversion. This connection is
an interpretation of their sequence. Augustine's Psalm exposition
explains the hope of the person crying from the depths; it does not
explain the Offertory's particular repetition. Fromage comments on the
actual offering and the petition that follows it.\footnote{Augustine,
\work{Enarrationes} 129.1--3; Fromage, vol. XI, pp. 512--513.}

\subsection{Secret: acceptance and converted desire}

The Secret asks God to accept his people's offerings and prayers and to
turn all their hearts toward him. Its two groups of desires are
distinguished by direction: earthly cravings and heavenly desires.
Lasance's ``enjoyments of heaven'' is a free historical rendering of
the latter. The Latin specifically asks for a change in what the
people desire.

The offering does not make conversion unnecessary. The people bringing
gifts ask for it within the same prayer. Fromage explains the passage
from earthly attachment toward heavenly good in his comment on this
Secret. Its application to the Gospel's person leaving possessions
behind will depend on the difference between using earthly goods and
treating them as the good that must never be lost.\footnote{Fromage,
vol. XI, p. 513.}

\subsection{Trinity Preface: the God who is thanked}

The prescribed Preface includes the dialogue, common thanksgiving,
confession of the Trinity and angelic conclusion. Its address to
the Father is followed by confession of the Son and Holy Spirit.
The prayer distinguishes persons while affirming one essence and
equal majesty. The historical English's ``without difference or
distinction'' must be read with its immediately following confession
of distinct persons: the equal divine glory is not a denial of
personal distinction.

Chrysostom observes in the Epistle that thanksgiving for the Father's
rescue does not exclude the Son's action. Redemption and forgiveness
belong to the Son in whom the people have been brought near.
Augustine likewise rejects an inference of inferior divinity from the
Son's or Spirit's being sent. These are doctrinal explanations that
illuminate the Preface, rather than commentaries on its wording.
The angels' unceasing praise continues into the received Sanctus
printed at the junction.\footnote{Chrysostom,
\work{Homilies on Colossians} 2, on 1:13--14; Augustine,
\work{De Trinitate} III, preface, §3.}

\subsection{Communion: the promise within Jesus' teaching on prayer}

Mark's account places the instruction about prayer after the disciples
have seen the withered fig tree, amid the journeys between Bethany
and Jerusalem and the Temple action. Jesus first requires faith in
God. He then speaks of the mountain, confident petition and forgiveness.
The Communion antiphon appoints the prayer promise from this instruction
in a distinct liturgical verbal form.\footnote{Mk 11:12--26.}

Bede's explanation concerns the adjacent sayings about faith and the
mountain, and forgiving while praying. His text does not separately
quote the appointed verse 24. He interprets powerful faith through
love and directs those who have not reached extraordinary perfection
to seek pardon while forgiving others. The second interpretation
develops why this immediate context matters to the antiphon's
apparently unlimited promise.\footnote{Bede,
\work{In Marci evangelium expositio} III, on Mk 11:22--23, 25.}

\subsection{Postcommunion: receiving and still needing cure}

The closing prayer speaks in the voice of people who have received the
mysteries. It asks for a cure of whatever is faulty in their minds.
The gift received is medicinal, and the prayer asks for its operation
in the recipients. Lasance's singular ``this sacrament'' represents
the Latin plural; its awkward ``may be restored'' concerns the person's
healing, not the preservation of a fault.

Augustine describes the sacramental food and drink in relation to the
Lord's Passion and spiritual health through God's invisible operation.
Fromage applies the actual Postcommunion to the cure of remaining
faults. Reception and further repentance can therefore belong
together. The prayer supplies no assurance that unrepented sin has
become harmless or that the communicant no longer needs to
persevere.\footnote{Augustine, \work{De Trinitate} III.4.10;
Fromage, vol. XI, p. 513; the prayer continues onto p. 514.}
